\exercisegroup

¶ 1) Hausdorff 局部凸空间的一个穷竭由一个筛 \(C = (C_n, p_n)_{n \geqslant 0}\) （GT, IX, § 6, No.~5, 定义 8）以及对每个 \(n \geqslant 0\) 、一个映射 \(\rho_n\) （从 \(C_n\) 到 E 的所有凸均衡子集所成的集中）给出，它具有下列性质：E1) E 是诸 \(\phi_0(c)\) （其中 \(c\) 取遍 \(C_0\) ）的并；

E2) 对每个 \(n\) 及每个 \(c \in \mathbf{C}_n\) ，\(\phi_n(c)\) 是诸 \(\phi_{n+1}(c')\) （其中 \(c'\) 取遍 \(p_n^{-1}(c)\) ）的并；E3) 对每个序列 \((c_k)_{k \geqslant 0}\) （满足 \(c_k \in \mathbf{C}_k\) 且 \(c_k = p_k(c_{k+1})\) ，对一切 \(k \geqslant 0\) ），存在一个序列 \((\phi_k)\) （由 \(>0\) 的数组成），使得对每个序列 \((x_k)\) （E 的点列，满足 \(x_k \in \phi_k(c_k)\) ）以及每个序列 \((\lambda_k)\) （实数的，满足 \(0 \leqslant \lambda_k \leqslant \rho_k\) ，对一切 \(k\) ），级数 \(\sum_{k=0}^{\infty} \lambda_k x_k\) 在 E 中收敛。

\begin{enumerate}
\def\labelenumi{\alph{enumi})}
\tightlist
\item
  在上述假设下，证明：若此外诸 \(\phi_n(c)\) （对 \(c \in \mathbf{C}_n\) ）还是闭的，则可以假设诸 \(\rho_k\) 已如此选取，使得 \(\sum_{k=m}^{\infty} \lambda_k x_k \in \phi_m(c_m)\) （对一切 \(m \geqslant 1\)）。

（取诸 \(\rho_{k}\) 使得 \(\sum_{k=0}^{\infty} \rho_{k} \leqslant 1\) ）。
\end{enumerate}

\begin{enumerate}
\def\labelenumi{\alph{enumi})}
\setcounter{enumi}{1}
\tightlist
\item
  假设给定一个筛 C 以及映到 E 的凸均衡子集所成之集中的映射的序列 \((\phi_{n})\) ，满足 E1)、E2) 以及下列条件：E3') 对每个序列 \((c_{k})_{k\geq0}\) （满足 \(c_{k}=p_{k}(c_{k+1})\) ，对一切 \(k\geq0\) ），存在一个序列 \((\mu_{k})\) （由 \textgreater0 的数组成），使得对每个序列 \((x_{k})\) （E 的点列，满足 \(x_{k}\in\phi_{k}(c_{k})\) ，对一切 k），点列 \((\mu_{k}x_{k})\) 含于 E 中一个凸的、有界均衡的且半完备的集中。
\end{enumerate}

证明此时条件 E3) 也成立（取 \(\rho_{k} = 2^{-k}\mu_{k}\) ）。

称局部凸 Hausdorff 空间为可穷竭的，如果存在 E 的一个穷竭。

\begin{enumerate}
\def\labelenumi{\arabic{enumi})}
\setcounter{enumi}{1}
\tightlist
\item
  设 \(\mathbf{E}\) 是一个局部凸空间且为 Baire 空间，\(\mathbf{F}\) 是一个局部凸可穷竭空间（习题 1），\((\mathbf{C}_n, p_n, \phi_n)\) 是 \(\mathbf{F}\) 的一个穷竭。
\end{enumerate}

\begin{enumerate}
\def\labelenumi{\alph{enumi})}
\item
  设 \(u\) 是从 \(\mathbf{E}\) 到 \(\mathbf{F}\) 中的一个线性映射，\(\mathbf{W}\) 是 \(\mathbf{F}\) 中一个凸的、均衡的且吸收的集。证明：存在一个序列 \((c_k)\) ，使得 \(c_k \in \mathbf{C}_k\) ，\(c_k = p_k(c_{k+1})\) （对一切 \(k \geqslant 0\) ），以及一个整数序列 \((m_k)\) （诸整数 \(> 0\) ），使得诸集 \(u^{-1}(\phi_k(c_k) \cap m_k\mathbf{W})\) （记作 \(\mathbf{M}_k\) ）中的每一个都不是 \(\mathbf{E}\) 中的贫集。证明：对每个 \(\varepsilon > 0\) ，存在一个数列 \((\nu_k)\) （诸数 \(> 0\) ），使得若点列 \((x_k)_{k \geqslant 1}\) （\(\mathbf{E}\) 的）满足 \(x_k \in \nu_k\mathbf{M}_k\) （对一切 \(k \geqslant 1\) ），则级数 \(\sum_{k=1}^{\infty} u(x_k)\) 在 \(\mathbf{F}\) 中收敛且其和属于 \(\varepsilon\mathbf{W}\) 。
\item
  此外假设 E 可度量化，且 u 在 \(E \times F\) 中的图像是闭的。证明：对每个 \(\varepsilon > 0\) ，有 \(\overline{u^{-1}(\mathbf{W})} \subset (1 + \varepsilon) u^{-1}(\mathbf{W})\) 。（注意：若 \((\mathrm{U}_{k})\) 是 E 中 0 的可数基本邻域系，则对每个 k 存在 E 中 0 的一个凸均衡邻域 \(V_{k}\) ，使得 \(V_{k} \subset U_{k} \cap v_{k} \overline{M}_{k}\) 。对每一点 \(a \in \overline{u^{-1}(\mathbf{W})}\) ，求一个序列 \((x_{k})_{k \geqslant 0}\) ，使得 \(x_{0} \in u^{-1}(\mathbf{W})\) ，\(x_{k} \in v_{k} M_{k}\) （对 \(k \geqslant 1\) ），且 \(a - \sum_{j=0}^{k} x_{j} \in V_{k}\) （对一切 \(k \geqslant 1\) ），然后应用 a)。
\item
  由 \(b\) ) 推出：若 \(E\) 是可度量化 Baire 空间，则从 \(E\) 到 \(F\) 中的每个图像为闭的线性映射都连续。
\end{enumerate}

\begin{enumerate}
\def\labelenumi{\arabic{enumi})}
\setcounter{enumi}{2}
\item
  证明：Fréchet 空间 \(\mathbf{E}\) 是可穷竭的（若 \((\mathbf{U}_k)\) 是构成 \(\mathbf{E}\) 中 0 的闭的凸均衡邻域的基本邻域系的一个递减序列，考虑诸集 \((m + 1)\mathrm{U}_k\) （其中 \(m\) 与 \(k\) 取遍 \(\mathbf{N}\) ）的有限交）。
\item
  \begin{enumerate}
  \def\labelenumii{\alph{enumii})}
  \tightlist
  \item
    可穷竭局部凸空间的每个闭子空间都是可穷竭的。
  \end{enumerate}
\end{enumerate}

\begin{enumerate}
\def\labelenumi{\alph{enumi})}
\setcounter{enumi}{1}
\tightlist
\item
  设 E 是一个可穷竭局部凸空间，\(u:E \rightarrow F\) 是从 E 到 F 上的一个连续线性满射映射。证明 F 是可穷竭的。特别地，E 对其闭子空间的每个商空间都是可穷竭的。在 E 上赋以比 E 的拓扑粗的 Hausdorff 局部凸拓扑所得到的每个空间都是可穷竭的。
\end{enumerate}

\begin{enumerate}
\def\labelenumi{\arabic{enumi})}
\setcounter{enumi}{4}
\tightlist
\item
  设 \((\mathrm{C}^{(m)})_{m\geqslant 0}\) 是筛 \(\mathrm{C}^{(m)} = (\mathrm{C}_n^{(m)},p_n^{(m)})_{n\geqslant 0}\) 的一个序列。对每个 \(n\geqslant 0\) ，令
\end{enumerate}

\[
\mathrm{D} _ {n} = \mathrm{C} _ {n} ^ {(0)} \times \mathrm{C} _ {n - 1} ^ {(1)} \times \dots \times \mathrm{C} _ {0} ^ {(n)} \times \prod_ {m = n + 1} ^ {\infty} \left\{a _ {m} \right\},
\]

其中 \(a_{m} = 0\) （对一切 \(m\geqslant 0\) ）；映射 \(p_n:\mathrm{D}_{n + 1}\to \mathrm{D}_n\) 取为等于

\[
p _ {n} ^ {(0)} \times p _ {n - 1} ^ {(1)} \times \dots \times p _ {0} ^ {(n)} \times q ^ {(n + 1)} \times \prod_ {m = n + 2} ^ {\infty} \mathrm{id} _ {m}
\]

其中 \(q^{(n + 1)}\) 是从 \(C_0^{(n + 1)}\) 到 \(\{0\}\) 上的唯一映射，\(\mathrm{id}_m\) 是 \(\{a_m\}\) 的恒等映射。那么 \((\mathrm{D}_n, p_n)\) 是一个筛。

设 \((\mathrm{E}^{(m)})_{m\geqslant 0}\) 是 Hausdorff 局部凸空间的一个序列；我们假设对每个 \(m\) 存在一个穷竭 \((\mathrm{C}_n^{(m)},p_n^{(m)},\phi_n^{(m)})_{n\geqslant 0}\) （\(\mathrm{E}^{(m)}\) 的）。考虑 Hausdorff 局部凸空间 \(\mathrm{E} = \prod \mathrm{E}^{(m)}\) ，且对每个 \(n\) ，令

\[
\phi_ {n} = \phi_ {n} ^ {(0)} \times \phi_ {n - 1} ^ {(1)} \times \dots \times \phi_ {0} ^ {(n)} \times \prod_ {m = n + 1} ^ {\infty} \psi_ {m}
\]

其中 \(\psi_{m}\) 是从 \(\{a_m\}\) 到 \(\mathrm{E}^{(m)}\) 的凸均衡子集所成之集中的映射，使得 \(\psi_{n}(a_{m}) = \mathrm{E}^{(m)}\) 。证明 \((\mathrm{D}_n,p_n,\phi_n)\) 是乘积空间 E 上的一个穷竭。

\begin{enumerate}
\def\labelenumi{\arabic{enumi})}
\setcounter{enumi}{5}
\tightlist
\item
  证明：递增的子空间序列 \(\mathbf{E}_n\) （向量空间 \(\mathbf{E}\) 的）的归纳极限（II, 第 31 页）（赋以拓扑 \(\mathcal{T}_n\) ，且 \(\mathbf{E}_n\) 赋以 \(\mathcal{T}_n\) 后可穷竭）若为 Hausdorff 的，则是一个可穷竭局部凸空间。
\end{enumerate}

